\chapter{Protocolo da revisão de escopo}
\label{ap:protocolo}

\section{Objetivo e tipo de revisão}

Esta pesquisa adota uma revisão de escopo (\textit{scoping review}) para
mapear a natureza, a extensão e as características das evidências
disponíveis sobre o uso de modelos de linguagem de grande porte
(\textit{Large Language Models} -- LLMs), inteligência artificial e
mecanismos de apoio à decisão no contexto do monitoramento e da gestão
da qualidade da água. A revisão busca identificar abordagens, tipos de
dados, aplicações, formas de interação com usuários e lacunas de pesquisa.
O relato do processo será organizado considerando as recomendações
aplicáveis do PRISMA-ScR \cite{tricco2018prismascr}.

A revisão de escopo tem como finalidade mapear a literatura, não presumir
antecipadamente a eficácia de uma intervenção nem substituir a avaliação
individual dos estudos. Por isso, as decisões de triagem por título e
metadados são preliminares: os trabalhos mantidos deverão passar por
leitura integral antes da inclusão definitiva na síntese.

\section{Questões de pesquisa}

As questões foram definidas para mapear a contribuição da literatura em quatro
camadas complementares: aplicações de IA generativa, integração com dados de
monitoramento, apoio à decisão e limitações das soluções propostas. A análise
será orientada por categorias de extração definidas a priori e refinadas durante
a leitura integral, preservando-se categorias emergentes quando os estudos
apresentarem aspectos não previstos no protocolo.

\subsection{QP1 -- Aplicações de LLMs e IA generativa}

\textbf{Questão.} Quais aplicações de LLMs, IA generativa, geração aumentada
por recuperação (RAG) ou agentes inteligentes têm sido propostas para o
monitoramento, a interpretação ou a gestão da qualidade da água e de sistemas
hídricos?

\textbf{Motivação.} Identificar os problemas hídricos que motivam o uso dessas
tecnologias e distinguir aplicações efetivas de propostas conceituais ou usos
genéricos de IA.

\textbf{Resposta esperada.} Um mapa de finalidades e tarefas, como interpretação
de dados, consulta em linguagem natural, diagnóstico, previsão, geração de
relatórios, recomendação e coordenação de agentes, indicando o contexto hídrico
quando informado.

\textbf{Estratégia de análise.} Extrair finalidade, contexto de aplicação,
tecnologia declarada e papel desempenhado pelo LLM/agente. Agrupar os estudos
por tarefa principal e registrar separadamente trabalhos conceituais, revisões
e implementações avaliadas empiricamente.

\subsection{QP2 -- Dados de sensores e monitoramento}

\textbf{Questão.} Quais tipos de dados de sensores, IoT ou redes de
monitoramento contínuo são utilizados, e como são tratados, analisados ou
interpretados?

\textbf{Motivação.} A utilidade de sistemas inteligentes para monitoramento
hídrico depende dos dados disponíveis, da frequência de coleta e da forma como
medições e informações contextuais são integradas.

\textbf{Resposta esperada.} Uma caracterização dos dados e fontes utilizados,
como parâmetros físico-químicos, séries temporais, dados operacionais,
observações remotas ou documentos técnicos, além das etapas de processamento
relatadas.

\textbf{Estratégia de análise.} Registrar sensores/fontes, parâmetros,
frequência ou temporalidade, pré-processamento, integração de dados e métodos
analíticos. Quando um aspecto não for informado, classificá-lo como não
reportado, sem inferir sua utilização.

\subsection{QP3 -- Mecanismos de apoio à decisão}

\textbf{Questão.} Quais mecanismos de apoio à decisão são descritos, incluindo
alertas, diagnóstico de anomalias, explicações, recomendações, visualizações e
interfaces humano--IA?

\textbf{Motivação.} O objetivo aplicado da pesquisa não se limita a identificar
modelos de IA: importa compreender como os resultados são transformados em
informações utilizáveis por operadores, especialistas ou gestores.

\textbf{Resposta esperada.} Uma comparação dos mecanismos de apoio à decisão,
dos usuários a que se destinam e da participação humana prevista no fluxo de
trabalho.

\textbf{Estratégia de análise.} Codificar o tipo de saída, o usuário-alvo, o
grau de automação, a possibilidade de intervenção humana e a forma de avaliação
do apoio à decisão, quando descritos no texto.

\subsection{QP4 -- Limitações e lacunas}

\textbf{Questão.} Quais limitações técnicas, metodológicas e de avaliação são
relatadas, e quais oportunidades de pesquisa permanecem em aberto?

\textbf{Motivação.} A identificação de limitações ajuda a evitar conclusões
baseadas apenas na novidade tecnológica e orienta a definição de requisitos
para arquiteturas inteligentes aplicáveis ao monitoramento da água.

\textbf{Resposta esperada.} Uma síntese das limitações relacionadas, por
exemplo, à qualidade e disponibilidade dos dados, confiabilidade das respostas,
validação empírica, explicabilidade, integração com sistemas existentes,
segurança e supervisão humana, apenas quando sustentadas pelos estudos.

\textbf{Estratégia de análise.} Extrair limitações e recomendações declaradas
pelos autores e agrupá-las por tema. Distinguir explicitamente limitações
relatadas pelos estudos de lacunas interpretadas na síntese desta revisão.

\section{Fontes de informação e datas das buscas}

As buscas foram realizadas em cinco fontes bibliográficas: Scopus,
Web of Science (WoS), ScienceDirect, IEEE Xplore e ACM Digital Library,
em 11 de outubro de 2026. Foram recuperados 135 registros brutos, assim
distribuídos: 20 na Scopus, 23 na WoS, 39 na ScienceDirect, 1 na IEEE
Xplore e 52 na ACM Digital Library. Esses valores correspondem aos
resultados registrados para a etapa de identificação da revisão.

\section{Estratégia de busca documentada}

As consultas utilizadas em cada fonte são apresentadas a seguir,
respeitando a sintaxe e os campos de busca de cada plataforma. Não foram
aplicados filtros adicionais nas consultas descritas.

\begin{description}
    \item[IEEE Xplore]\mbox{}\\
    {\small\ttfamily\raggedright
    "water quality monitoring" AND "decision support" AND "sensor data" AND "large language models"\par}
    Sem filtros.

    \item[ScienceDirect]\mbox{}\\
    {\small\ttfamily\raggedright
    "water quality monitoring" AND "decision support" AND "sensor data" AND "large language models"\par}
    Sem filtros.

    \item[Web of Science]\mbox{}\\
    {\small\ttfamily\raggedright
    "water quality monitoring" (Topic) AND "decision support" (Topic) AND "sensor data" (Topic) AND "large language models" (Topic)\par}
    Sem filtros.

    \item[Scopus]\mbox{}\\
    {\small\ttfamily\raggedright
    "water quality monitoring" AND "decision support" AND "sensor data" AND "large language models"\par}
    Campo: \textit{All fields}.

    \item[ACM Digital Library]\mbox{}\\
    {\small\ttfamily\raggedright
    [All: "water quality monitoring"] AND [All: "decision support"] AND [All: "sensor data"] AND [All: "large language models"]\par}
    Sem filtros.
\end{description}

As expressões combinam quatro conceitos escolhidos por corresponderem ao
problema aplicado (monitoramento da qualidade da água), à finalidade decisória
(apoio à decisão), à fonte de dados de interesse (dados de sensores) e à família
de tecnologia investigada (modelos de linguagem de grande porte). O operador
AND foi usado para exigir a presença dos quatro conceitos na recuperação. Essa
opção aumenta a especificidade, mas pode reduzir a sensibilidade e deixar de
recuperar trabalhos que usem sinônimos ou não mencionem todos os conceitos nos
campos pesquisados. A sintaxe e os campossão similares mas variam por
plataforma e portanto são documentados individualmente para favorecer a transparência e
a reprodutibilidade.
\section{Critérios de elegibilidade}

Os critérios abaixo constituem a regra de triagem proposta para esta
versão e deverão ser aplicados consistentemente na leitura dos textos completos.

\subsection{Critérios de inclusão}

\begin{itemize}
    \item Publicações individuais potencialmente relacionadas a pelo menos
    uma questão de pesquisa.
    \item Artigos de periódicos e artigos individuais de conferências, incluindo
    estudos primários e revisões pertinentes, com publicação formal e revisão
    por pares verificável.
    \item Estudos sobre monitoramento ou qualidade da água; dados de
    sensores, IoT ou monitoramento contínuo; LLMs, IA generativa, RAG
    ou agentes; ou mecanismos de apoio à decisão em sistemas hídricos.
    \item Publicações com informação suficiente para uma possível
    contribuição. Primeiramente sendo feita uma consulta ao resumo e, quando elegíveis, ao texto
    completo.
\end{itemize}

\subsection{Critérios de exclusão}

\begin{itemize}
    \item Registros duplicados, preservando-se uma ocorrência por publicação.
    \item Registros de patentes, conjuntos de dados, índices de livros ou
    volumes inteiros de anais, quando não correspondam a um estudo
    individual.
    \item Livros completos e capítulos de livros, pois não serão considerados
    parte elegível desta pesquisa. A unidade de análise será o artigo
    individual publicado em periódico ou conferência com revisão por pares
    verificável.
    \item Publicações cujo conteúdo não tenha relação suficientemente
    direta com as questões de pesquisa.
    \item Manuscritos não publicados, preprints ou outros registros sem
    publicação formal verificável até a data da busca, caso não atendam ao
    recorte de literatura publicada adotado nesta revisão.
    \item Publicações para as quais, após verificação, não seja possível
    confirmar revisão por pares, conforme o recorte adotado para o corpus.
    A ausência inicial dessa informação será tratada como pendência de
    verificação, não como prova automática de ausência de revisão por pares.
    \item Texto completo indisponível somente após tentativa documentada
    de obtenção.
\end{itemize}

A ausência de metadados suficientes sobre a revisão por pares será
registrada como limitação da triagem bibliográfica. Como não haverá leitura
integral nesta etapa, os resultados serão descritos como seleção preliminar
por título e metadados, sem alegar confirmação de elegibilidade em texto completo.

\section{Procedimento de seleção}

A seleção foi realizada por identificação dos registros nas cinco fontes,
deduplicação e triagem bibliográfica por tipo documental, título e metadados.
Dos 30 registros submetidos à avaliação complementar de pertinência, 13 foram
excluídos por incompatibilidade temática. Dos 17 restantes, dois capítulos de
livro foram excluídos por não atenderem ao recorte documental adotado. Os 15
registros restantes compõem o conjunto selecionado para a análise bibliográfica
preliminar apresentada no Apêndice B. Não foi realizada leitura integral nesta
etapa, portanto, a seleção não deve ser interpretada como validação de todos
os detalhes metodológicos ou de elegibilidade a partir do texto completo.

\section{Estratégia de snowballing e avaliação da qualidade}

Foi prevista a consulta às listas de referências dos estudos mantidos após a
triagem bibliográfica, como procedimento complementar para identificar
publicações potencialmente relevantes que não tivessem sido recuperadas pelas
buscas nas bases. A execução e o resultado desse procedimento são relatados no
Apêndice B.

Não foi definido um instrumento formal de avaliação da qualidade metodológica
para esta etapa, pois o objetivo imediato é mapear a literatura e descrever suas
características. Isso não significa que a qualidade dos estudos tenha sido
considerada equivalente: limitações de método, validação e evidência empírica
serão extraídas e descritas quando a leitura integral for realizada.

\section{Extração e síntese dos dados}

Para cada estudo incluído após leitura integral, serão extraídos, quando
disponíveis: referência bibliográfica, ano, veículo de publicação,
contexto hídrico, objetivo, tipos de dados, sensores ou parâmetros
monitorados, método de processamento, papel dos LLMs/agentes, mecanismo
de apoio à decisão, forma de avaliação, resultados, limitações e lacunas.
A síntese combinará estatísticas descritivas e categorização qualitativa
orientada pelas questões de pesquisa.

\section{Limitações do registro metodológico}

Quando o ano de fascículo for posterior à data da busca, será considerada
a data de disponibilização eletrônica da publicação. Assim, a indicação
de um ano futuro no fascículo não implica, por si só, que o registro ainda
não estivesse publicado on-line na data da busca.
